%% toc: P15 173 173.9,173.10,173.11,173.12
\chapter{Computed and Simulated Values in Approximate Conclusions}
\label{sec:p15-w02b-p15-074}

\Cref{sec:p15-approximate-quantitative-inference} treats an approximation
as a stated relation between an exact model value and a simpler expression.
In use, the simpler expression, or the exact model itself, is evaluated by
a program, and the number the program returns is a further object: a
computed value. This chapter states how a computed or simulated value enters
a sign, ordering, threshold or regime argument. The rule is bookkeeping.
Each error term refers to one explicit difference between two named
quantities, each term is bounded on its own premises, and the bounds are
combined only afterwards. Arithmetic error, discretization error,
approximation error, input uncertainty and the discrepancy between a model
and its target are different differences. A bound on one of them says
nothing about another, and making one of them small leaves the others in
the conclusion.

\emph{Inputs.} A target quantity and a strict comparison to be decided;
an exact model and, where used, an approximant with a remainder bound on a
stated domain; a recorded input with an enclosure of the input the model
should receive; a computed value with a bound on its arithmetic and, for a
discretized calculation, discretization error; for a sampled simulation, a
number of independent runs. \emph{Conclusions.} An interval certified to
contain the target, a certified strict comparison or a justified abstention,
a certified regime condition, and, for a sampled simulation, a decision
rule with a stated error probability. When a comparison is not certified,
the analysis also identifies which term prevents it.

\emph{Dependencies.} The approximation relations of
\cref{def:p15-approximation-relations}, the interval certificates of
\cref{thm:p15-approximation-certificates} and the remainder bounds of
\cref{thm:p15-approximation-rational} are used as stated there. From
\cref{part:p12}, the chapter uses the finite radix format and nearest
rounding of \cref{def:p12-numerical-format,def:p12-numerical-rounding},
the half-ulp bound of \cref{thm:p12-numerical-nearest}, the sign and
ordering certificates of \cref{thm:p12-numerical-certification}, the
separation of input uncertainty from arithmetic error in
\cref{thm:p12-numerical-data-budget}, the scalar error chain of
\cref{def:p12-simulation-error-chain,thm:p12-simulation-error-chain}, the
distinctions of \cref{def:p12-simulation-vv} and the Euler results
\cref{thm:p12-simulation-decay-euler,thm:p12-simulation-arithmetic-propagation}.
The sampling result uses \cref{thm:p15-expectation-bounds},
\cref{thm:p15-independent-blocks} and \cref{thm:p15-independent-moments}.
Within the chapter, \cref{thm:p15-w02b-p15-074-ledger} rests only on the
triangle inequality; \cref{thm:p15-w02b-p15-074-regime} applies it to a
regime parameter; the two worked theorems
\cref{thm:p15-w02b-p15-074-four-thresholds,thm:p15-w02b-p15-074-euler}
apply both; and \cref{thm:p15-w02b-p15-074-sampled,thm:p15-w02b-p15-074-tie}
are independent of the deterministic results except for the rounding model.

\section{The ledger of separate differences}
\label{sec:p15-w02b-p15-074-ledger}

\begin{hbdefinition}{def:p15-w02b-p15-074-ledger}{An error ledger for a computed value}
Let \(Y\) be a real target quantity in stated units. Let \(F\) be a real
model function on an input set \(S\), and let \(x_*\in S\) denote the
unknown input that the model should receive, so that \(F(x_*)\) is the
model value. Let \(A\) be an approximant, a real function on \(S\) intended
to replace \(F\); when the model itself is evaluated, \(A=F\). Let
\(\widehat x\in S\) be the recorded input, \(B\subseteq S\) an
\emph{input enclosure} containing both \(x_*\) and \(\widehat x\), and
\(\widehat A\) the number returned by a computation intended to evaluate
\(A(\widehat x)\). An \emph{error ledger} consists of nonnegative numbers
\(E_{\rm model}\), \(E_{\rm input}\), \(E_{\rm comp}\) and real numbers
\(\ell_{\rm appr}\leq r_{\rm appr}\) such that
\[
\begin{aligned}
 &|Y-F(x_*)|\leq E_{\rm model},&&
 \ell_{\rm appr}\leq F(x)-A(x)\leq r_{\rm appr}\ \ (x\in B),\\
 &|A(x)-A(\widehat x)|\leq E_{\rm input}\ \ (x\in B),&&
 |A(\widehat x)-\widehat A|\leq E_{\rm comp}.
\end{aligned}
\]
In the symmetric case \(\ell_{\rm appr}=-E_{\rm appr}\) and
\(r_{\rm appr}=E_{\rm appr}\); when \(A=F\), both are zero. If the
computation is a discretized calculation whose exact-real result
\(A_h(\widehat x)\) satisfies \(|A(\widehat x)-A_h(\widehat x)|\leq
E_{\rm disc}\) and \(|A_h(\widehat x)-\widehat A|\leq E_{\rm arith}\),
then \(E_{\rm comp}=E_{\rm disc}+E_{\rm arith}\) is admissible.

Each entry is a separate claim with its own evidence. The model entry
concerns correspondence between the model and its target; the
approximation entry is a remainder theorem holding on all of \(B\), not
only at \(\widehat x\); the input entry combines an enclosure of \(x_*\)
with a sensitivity bound for \(A\); the arithmetic entry is a rounding
analysis of the operations actually executed, under the range conditions
of \cref{def:p12-numerical-format}; the discretization entry is a
convergence or error theorem for the discrete method. A bound that has not
been established is recorded as unavailable. It is not set to zero.
\index{error ledger}\index{arithmetic error!separated from model error}
\end{hbdefinition}

The ledger arranges the four differences in a telescoping order:
\[
 Y-\widehat A=\bigl(Y-F(x_*)\bigr)+\bigl(F(x_*)-A(x_*)\bigr)
 +\bigl(A(x_*)-A(\widehat x)\bigr)+\bigl(A(\widehat x)-\widehat A\bigr).
\]
It differs from the chain of \cref{def:p12-simulation-error-chain} in two
respects: it inserts an approximant with a possibly one-sided remainder,
and it applies the input perturbation to the approximant, the function
actually evaluated. Either order is legitimate provided the terms refer to
explicit differences and none is counted twice.

\begin{hbtheorem}{thm:p15-w02b-p15-074-ledger}{proved=annex:p15-w02b-p15-074-ledger}{A ledger as a threshold certificate}
Under the premises of \cref{def:p15-w02b-p15-074-ledger}, put
\(E=E_{\rm model}+E_{\rm input}+E_{\rm comp}\). Then
\[
 \widehat A+\ell_{\rm appr}-E\;\leq\;Y\;\leq\;\widehat A+r_{\rm appr}+E.
\]
\begin{enumerate}
\item For a real threshold \(c\), the inequality
\(\widehat A+\ell_{\rm appr}-E>c\) certifies \(Y>c\), and
\(\widehat A+r_{\rm appr}+E<c\) certifies \(Y<c\).
\item For two targets \(Y_1,Y_2\) with ledgers indexed accordingly,
\(\widehat A_1+r_{1}+E_1<\widehat A_2+\ell_{2}-E_2\) certifies
\(Y_1<Y_2\), where \(r_1\) and \(\ell_2\) are the respective
approximation endpoints.
\item Regard the four premises as constraints on the five real numbers
\(Y\), \(F(x_*)\), \(A(x_*)\), \(A(\widehat x)\), \(\widehat A\), with
\(\widehat A\) fixed. The set of values of \(Y\) compatible with them is
exactly the displayed closed interval. Hence, if \(c\) lies in that
interval, the premises alone certify neither \(Y>c\) nor \(Y<c\).
\item In the symmetric case, if at least two of \(E_{\rm model}\),
\(E_{\rm appr}\), \(E_{\rm input}\), \(E_{\rm comp}\) are positive, the
root-sum-square of the four is strictly smaller than their sum and is
exceeded by \(|Y-\widehat A|\) for some assignment satisfying the premises.
It is therefore not a valid worst-case bound under these premises.
\end{enumerate}
\end{hbtheorem}

When \(A=F\) and \(E_{\rm model}=E_{\rm input}=0\), the first item reduces
to the sign certificate of \cref{thm:p12-numerical-certification}; when
\(A=F\) and the computation is discretized, the inequality reduces to
\cref{thm:p12-simulation-error-chain}. A root-sum-square combination can be
justified as a statement about standard deviations under a probability
model with uncorrelated terms; it is then a different kind of claim, as
\cref{sec:p15-measurement-uncertainty} explains, and it does not certify a
strict comparison by itself.

A computed value is also used to decide which approximation applies at all.
The terminal-speed check of \cref{def:p15-approximation-drag} substitutes a
computed speed into a Reynolds number, and the remainder bound of
\cref{thm:p15-approximation-rational} holds only for inputs up to one half.
Such a regime condition is a threshold statement about the input the model
should receive, so it is decided by the same certificate applied to the
regime parameter.

\begin{hbcorollary}{thm:p15-w02b-p15-074-regime}{proved=annex:p15-w02b-p15-074-regime}{A regime is certified on the enclosure, not at the computed value}
Let \(\varphi\) be a real function on \(S\), \(\rho\) a real bound, and
\(D=\{x\in S:\varphi(x)\leq\rho\}\) the domain on which a remainder bound
has been proved. Suppose \(|\varphi(x)-\varphi(\widehat x)|\leq H_{\rm input}\)
for every \(x\in B\), and a computed value \(\widehat\varphi\) satisfies
\(|\varphi(\widehat x)-\widehat\varphi|\leq H_{\rm arith}\). If
\[
 \widehat\varphi+H_{\rm input}+H_{\rm arith}\leq\rho,
\]
then \(B\subseteq D\), and the remainder bound may serve as the
approximation entry of the ledger on \(B\). The weaker premise
\(\widehat\varphi\leq\rho\) does not imply \(B\subseteq D\), even when
\(\widehat\varphi=\varphi(\widehat x)\) exactly.
\end{hbcorollary}

The corollary needs a numerical bound \(\rho\). A regime described only as
``small'' or ``large'' supplies no threshold to certify against. In the
falling-sphere example the computed Reynolds number is about \(0.016\).
Against a stated bound \(\rho=1\), the corollary needs only that the input
and arithmetic entries for the Reynolds number sum to at most
\(1-\widehat\varphi\), roughly \(0.98\). The certificate still requires
those entries to be bounded and the bound \(\rho\) to be stated.

\section{Four thresholds and the term that blocks each}
\label{sec:p15-w02b-p15-074-four}

The following calculation keeps every entry of the ledger explicit. Its
arithmetic uses the decimal format of precision three that Goldberg uses
throughout his discussion of ulps and relative error
\cite[Section D.3.2]{GoldbergOracle8193693}, with the nearest rounding
\(R\) of \cref{def:p12-numerical-rounding}. A small precision makes each
rounding visible; nothing in the argument depends on that choice.

\begin{hbtheorem}{thm:p15-w02b-p15-074-four-thresholds}{proved=annex:p15-w02b-p15-074-four-thresholds}{One computed value, four thresholds}
Let \(x_*\) be a dimensionless ratio, \(F(x)=1/(1-x)\) and
\(A(x)=1+x\). Suppose the recorded value is \(\widehat x=0.123\), with
\(|x_*-\widehat x|\leq0.002\), so \(B=[0.121,0.125]\). Let
\(\widehat A=R(1+0.123)\), one nearest-rounded addition in decimal
precision three. Take the target to be \(Y=F(x_*)\), so
\(E_{\rm model}=0\), except in item 6.
\begin{enumerate}
\item \(\widehat A=1.12\). A ledger holds with \(E_{\rm comp}=0.005\),
\(E_{\rm input}=0.002\), \(\ell_{\rm appr}=0\) and
\(r_{\rm appr}=1/32\), and \(B\subseteq[0,1/2]\). The certified interval
is \([1.113,\,1.15825]\).
\item The ledger certifies \(Y>1.10\).
\item The interval contains \(1.12\), so this ledger certifies neither
strict comparison with \(1.12\). Replacing \(\widehat A\) by the exact
value \(A(\widehat x)=1.123\), so that \(E_{\rm comp}=0\), gives the interval
\([1.121,\,1.15625]\) and certifies \(Y>1.12\).
\item With \(E_{\rm comp}=0\) the interval still contains \(1.13\).
Evaluating \(F\) itself on \(B\) gives
\(F(B)=[1000/879,\,8/7]\), whose lower endpoint exceeds \(1.137\);
this certifies \(Y>1.13\).
\item The value \(1.14\) lies in \(F(B)\). The inputs \(0.121\) and
\(0.125\) both satisfy the input premise, and \(F\) takes values below and
above \(1.14\) at them. No reduction of the arithmetic or approximation
entries decides the comparison of \(Y\) with \(1.14\).
\item If instead \(Y\) is a target quantity with
\(|Y-F(x_*)|\leq E_{\rm model}\), the certificate in item 2 persists
exactly when \(E_{\rm model}<0.013\), and that in item 4 exactly when
\(E_{\rm model}<673/87900\), approximately \(0.00766\).
\end{enumerate}
\end{hbtheorem}

\begin{hbexample}{def:p15-w02b-p15-074-four-reading}{Reading the four outcomes}
In item 3 the printed value \(1.12\) equals the threshold. That
coincidence is not evidence for equality or for either strict inequality:
a computed value supports no comparison until its ledger is attached. The rounding error actually committed is \(0.003\),
smaller than the budget \(0.005\); a pessimistic budget yields abstention,
never a false certificate.

The four thresholds are blocked by different terms. At \(1.10\) every
generic bound leaves margin. At \(1.12\) only the arithmetic entry blocks,
and exact evaluation of the same approximant suffices, because the input
entry alone leaves a margin of \(0.001\). At \(1.13\) the first-order
approximant fails even when evaluated exactly: more digits cannot repair
an approximation entry. The exact model, or the sharper lower remainder
\(x^2/(1-x)\geq x^2\), which gives \(1.121+0.121^2=1.135641\), supplies
the missing margin. At \(1.14\) only better information about \(x_*\), a
narrower enclosure, can help. Item 6 shows the remaining entry: a
certificate about \(F(x_*)\) transfers to a measured target only through a
model-to-target bound, and a small one can already undo it.

If \(\widehat x\) had itself been computed, for instance as a rounded
quotient of two recorded quantities, that rounding error would enlarge the
enclosure \(B\) before the regime check of
\cref{thm:p15-w02b-p15-074-regime} and before every later entry.
The same separation is the point of the discriminant example in
\cref{thm:p12-numerical-data-budget,def:p12-numerical-certification-cases}:
an exact final subtraction can expose rounding committed earlier, the
catastrophic cancellation that Goldberg distinguishes from the benign case
of subtracting exactly known operands
\cite[Section D.3.4]{GoldbergOracle8193693}.
\end{hbexample}

\section{A simulated value in a threshold argument}
\label{sec:p15-w02b-p15-074-simulation}

A deterministic simulation returns a computed value of a discretized
calculation. Its ledger therefore has both a discretization and an
arithmetic entry, and the regime conditions of the discretization theorem
are among the premises. The NASA simulation document lists finite-precision
arithmetic beside discretization, incomplete iterative convergence and
numerical approximation as numerical errors whose estimates should be
recorded for a computational model's results, and it places solution
verification at each use of a model
\cite[Sections 4.2.3--4.2.4]{NASA7009B2024}. Its uncertainty
characterization keeps input uncertainty, numerical error in the
implementation and propagation to outputs as distinct aspects
\cite[Section E.4.4]{NASA7009B2024}. The following theorem shows those
distinctions at work on the exact decay model of
\cref{def:p12-simulation-decay}.

\begin{hbtheorem}{thm:p15-w02b-p15-074-euler}{proved=annex:p15-w02b-p15-074-euler}{Thresholds from an Euler calculation}
Consider \(q'=-q\), \(q(0)=1\), whose model value at time one is
\(q(1)\), and treat that value as unknown. Let \(v_N=(1-1/N)^N\) be the
exact-real Euler iterate with \(N\) steps of size \(h=1/N\), and let
\(\widehat v_{j+1}=R(0.95\,\widehat v_j)\), \(\widehat v_0=1\), be the
recurrence for \(N=20\) computed in decimal precision three.
\begin{enumerate}
\item \(0.35848<v_{20}\leq q(1)\leq\tfrac{38}{37}v_{20}<0.36818\).
\item \(\widehat v_{20}=0.358\), and
\(|\widehat v_{20}-v_{20}|\leq0.01\,(1-v_{20})<0.0065\).
\item The computed value certifies \(q(1)>0.35\). With this arithmetic
budget it does not certify \(q(1)>0.355\), whereas the exact-real iterate
\(v_{20}\) does.
\item At \(N=20\), even exact-real iteration certifies neither strict
comparison of \(q(1)\) with \(0.365\). With \(N=100\), \(v_{100}>0.365\)
certifies \(q(1)>0.365\).
\end{enumerate}
\end{hbtheorem}

In item 4 the comparison \(q(1)>0.365\) is true, since
\(q(1)=e^{-1}\approx0.3679\); the abstention at \(N=20\) is a statement
about what the discretization bound establishes, not about the value. The
step size satisfies the regime condition \(0<\lambda h<1\) of
\cref{thm:p12-simulation-decay-euler} with \(\lambda=1\) known exactly; with an uncertain
rate, that condition would itself require the enclosure argument of
\cref{thm:p15-w02b-p15-074-regime}. Every item concerns the mathematical
model. A conclusion about a physical decay needs a model-to-target entry,
and \cref{thm:p12-simulation-regime-change} exhibits two trajectories that
agree exactly up to time one and lie on opposite sides of \(0.1\) at time
two. Refining the step or the arithmetic converges to the model's value,
whichever trajectory the target follows.

\section{Sampled simulations near a threshold}
\label{sec:p15-w02b-p15-074-sampled}

A stochastic simulation often estimates the probability that a model
output reaches a threshold, by the frequency of runs in which the computed
output does so. Two further differences then enter: the computed indicator
can differ from the exact one, and the sampled frequency differs from the
probability it estimates. The first is controlled by arithmetic error and by
how much probability lies near the threshold; the second by the number of
runs.

\begin{hbtheorem}{thm:p15-w02b-p15-074-sampled}{proved=annex:p15-w02b-p15-074-sampled}{A sampled frequency with computed indicators}
Let \(\xi\) be a random element of a measurable space \(S\), let
\(G,\widehat G:S\to\mathbb R\) be measurable, let \(a\) be real, and
suppose \(|\widehat G(s)-G(s)|\leq E\) for every \(s\in S\). Put
\[
 p=\mathbb P(G(\xi)\geq a),\qquad
 \widetilde p=\mathbb P(\widehat G(\xi)\geq a),\qquad
 \omega(E)=\mathbb P(|G(\xi)-a|\leq E).
\]
\begin{enumerate}
\item \(|\widetilde p-p|\leq\omega(E)\).
\item If \(\xi_1,\ldots,\xi_n\) are independent with the law of \(\xi\)
and \(\widehat p_n\) is the fraction of indices \(i\) with
\(\widehat G(\xi_i)\geq a\), then for every \(t>0\),
\[
 \mathbb P\bigl(|\widehat p_n-\widetilde p|\geq t\bigr)\leq\frac1{4nt^2}.
\]
\item For a probability threshold \(p_0\), if \(p\leq p_0\) then
\(\mathbb P\bigl(\widehat p_n>p_0+t+\omega(E)\bigr)\leq1/(4nt^2)\);
if \(p\geq p_0\) then
\(\mathbb P\bigl(\widehat p_n<p_0-t-\omega(E)\bigr)\leq1/(4nt^2)\).
\end{enumerate}
\end{hbtheorem}

Item 3 is a decision rule with a stated error probability, not a
worst-case certificate: declaring \(p>p_0\) when
\(\widehat p_n>p_0+t+\omega(E)\) errs with probability at most
\(1/(4nt^2)\) whenever \(p\leq p_0\). The sampling bound concerns
\(\widetilde p\), the probability for the computed indicator. More runs
shrink the sampling margin \(t\) at fixed error probability; they do
not shrink \(\omega(E)\). Bounding \(\omega(E)\) needs information about
the law of \(G(\xi)\) near \(a\); a density bounded by \(K\) on
\([a-E,a+E]\), for instance, gives \(\omega(E)\leq2KE\). An atom at the
threshold gives no such bound, as the next result shows.

\begin{hbtheorem}{thm:p15-w02b-p15-074-tie}{proved=annex:p15-w02b-p15-074-tie}{A tie at the threshold, decided by rounding}
Let \(\xi\) be uniform on \(\{1,2,3\}\), \(G(s)=(s/3)\cdot3\) and
\(a=1\), so that \(p=1\). Let \(\widehat G(s)=R\bigl(R(s/3)\cdot3\bigr)\)
in decimal precision three. Then \(\widehat G(1)=0.999\),
\(\widehat G(2)=2\) and \(\widehat G(3)=3\); the uniform arithmetic
bound \(E=0.001\) holds, \(\widetilde p=2/3\) and \(\omega(E)=1/3\),
so item 1 of \cref{thm:p15-w02b-p15-074-sampled} holds with equality.
For the threshold \(p_0=0.9\) and any \(t>0\), the rule of item 3
never declares \(p>0.9\), and it declares \(p<0.9\) with probability at
most \(1/(4nt^2)\). If \(\omega(E)\) is omitted from the rule and
\(0<t<7/30\), it declares the false \(p<0.9\) with probability at least
\(1-1/\bigl(4n(7/30-t)^2\bigr)\).
\end{hbtheorem}

The last bound tends to one as \(n\) grows: more runs make the error
from ignoring \(\omega(E)\) more likely, not less, because the
frequencies concentrate near \(2/3\). The arithmetic entry here is
one unit in the third digit, yet it moves a third of the probability,
because that probability sits exactly at the threshold.

\section{Limits of the ledger and a procedure}
\label{sec:p15-w02b-p15-074-limits}

\begin{hbwarning}{def:p15-w02b-p15-074-limits}{What the ledger does not supply}
\begin{enumerate}
\item A standard uncertainty is not an error bound
(\cref{def:p15-measurement-error}). Using \(ku\) as \(E_{\rm input}\)
turns a dispersion statement into a stipulated enclosure. The resulting
certificate holds conditionally on that enclosure containing \(x_*\), with
whatever probability the enclosure's own interpretation supports
(\cref{thm:p15-measurement-coverage-models}).
\item Generic rounding bounds hold under the range conditions of
\cref{def:p12-numerical-format}. Overflow, underflow, exceptional values
and flush-to-zero behavior require a separate analysis.
\item Agreement between two runs, two compilers or two machines is not an
arithmetic-error bound: the same rounding errors can repeat. Agreement of
leading digits as the step size shrinks suggests convergence; it is a
discretization bound only when a theorem or a stated error estimate
connects it to the exact-real result.
\item The probability in \cref{thm:p15-w02b-p15-074-sampled} assumes
independent draws with the law of \(\xi\). A pseudo-random generator
produces a deterministic sequence; treating it as such draws is a modeling
premise, not a consequence of the theorem.
\item Solution verification bounds the computation entries. It supplies no
model-to-target entry; that requires validation evidence in the sense of
\cref{def:p12-simulation-vv}.
\item An asymptotic relation \(f\sim g\) or \(f=g+O(h)\)
(\cref{def:p15-approximation-relations}) supplies no approximation entry
until its constants and its domain are stated. A computed value of \(g\)
then has no certified relation to \(f\) at a finite input.
\end{enumerate}
\end{hbwarning}

The worst-case ledger can be pessimistic, as item 3 of
\cref{thm:p15-w02b-p15-074-four-thresholds} shows. Sharper statements
need sharper premises: an exact evaluation, a better remainder, a narrower
enclosure, or a probability model for the errors. A probabilistic rounding
model makes a different claim from the bounds used here and needs its own
justification.

\begin{hbmethod}{def:p15-w02b-p15-074-procedure}{Entering a computed or simulated value into an approximate conclusion}
\begin{enumerate}
\item State the target, its units and the strict comparison to be decided.
\item Name the model, the approximant and the relation between them in the
sense of \cref{def:p15-approximation-relations}; obtain a remainder bound
with numerical constants on a stated domain.
\item Record the inputs and an enclosure of the inputs the model should
receive, including rounding committed in deriving the recorded inputs.
\item Certify every regime condition on the enclosure by
\cref{thm:p15-w02b-p15-074-regime}.
\item Bound the arithmetic error of the executed operations and, for a
discretized calculation, the discretization error, separately.
\item State the model-to-target entry, or record it as unavailable and
restrict the conclusion to the model value.
\item Combine the entries by \cref{thm:p15-w02b-p15-074-ledger} and apply
the certificate. If the interval contains the threshold, compare each
entry with the margin to find the one that blocks, and improve that one,
or report the comparison as undecided.
\item For a sampled simulation, add the indicator band of
\cref{thm:p15-w02b-p15-074-sampled} and report the decision with its
error probability.
\end{enumerate}
This extends \cref{def:p15-approximation-procedure} by the computation and
regime entries, and \cref{def:p12-numerical-inference-procedure} by the
approximation entry.
\index{threshold decision!computed value}\index{simulation!threshold decision}
\end{hbmethod}

\section{Relationships and sources}
\label{sec:p15-w02b-p15-074-relations}

The certificates of \cref{thm:p15-approximation-certificates} and
\cref{thm:p12-numerical-certification} are the special cases of
\cref{thm:p15-w02b-p15-074-ledger} with a single entry. The threshold
reversals of \cref{thm:p15-approximation-rational} and the nonuniform
convergence of \cref{thm:p15-approximation-nonuniform} concern the
approximation entry; the corollary on regimes adds that an input near a
transition must be enclosed, not merely computed, before a remainder bound
applies. The range estimate of \cref{sec:p15-w02b-p15-073} computes the
exact image of an input box, which is the input entry made exact for a
monotone model. The probabilistic treatment of input uncertainty is in
\cref{sec:p15-measurement-uncertainty}. In \cref{part:p12},
\cref{sec:p12-numerical-error-and-inference} develops the arithmetic
entry, and \cref{sec:p12-simulation-results-as-evidence} develops the
discretization, verification, validation and transfer questions,
including \cref{thm:p12-simulation-finite-comparisons} on why finitely many
agreements give no uniform prediction bound.

Sources: the decimal precision-three format and the half-ulp and
relative-error measures follow \cite[Section D.3.2]{GoldbergOracle8193693};
the distinction between catastrophic and benign cancellation follows
\cite[Section D.3.4]{GoldbergOracle8193693}. The classification of numerical
errors and the place of solution verification follow
\cite[Sections 4.2.3--4.2.4 and E.4.4]{NASA7009B2024}. The ledger, its
certificate and sharpness, the regime corollary, the worked thresholds, the
Euler threshold calculation and the two sampling results are derivations
of this handbook.
